% Abstract (BRM limit 250 words). Its numbers carry the receipts of the body sections.
A language model prompted with a demographic persona can give answers that each look plausible while the simulated population sits far from the real one.
Researchers still use such samples to estimate scores, compare groups and score scales for populations they cannot easily survey.
We propose the validity ladder, a set of checks released as code, in which each use of a simulated sample is tested separately against human data with its own statistic and pass rule.
A sample earns only the claims whose checks it passes, and researchers can set each pass threshold to the risk of their application.
Before any answers are generated, a certifiability step uses the human data alone to show which group differences a study could confirm.
Real survey respondents run through the same checks rarely fail, and errors planted at known sizes are caught by the check built for each.
Across ten datasets released by other authors, including silicon samples and digital twins, samples drawn from the datasets' own human answers rarely fail where the human data are precise.
None of 725 readings of published simulators in nine of them reproduces the group differences, and for over half the human data were too thin to confirm any simulator.
In 28,800 PHQ-8 depression assessments from four models, persona averages stabilize after a few answers, yet every model fails every check above that.
The ladder gives researchers a checkable list of the claims a simulated sample supports and the human data needed before generating it.
